% =====================================================================
% Central regularisation of the Igusa reciprocal for K3 x E.
% =====================================================================

\providecommand{\kBKM}{\kappa_{\mathrm{BKM}}}
\providecommand{\kch}{\kappa_{\mathrm{ch}}}
\providecommand{\kcat}{\kappa_{\mathrm{cat}}}
\providecommand{\kfib}{\kappa_{\mathrm{fiber}}}
\providecommand{\kHeis}{\kappa_{\mathrm{ch}}^{\mathrm{Heis}}}
\providecommand{\Sp}{\mathrm{Sp}}
\providecommand{\SL}{\mathrm{SL}}
\providecommand{\DT}{\mathrm{DT}}
\providecommand{\GW}{\mathrm{GW}}
\providecommand{\MNOP}{\mathrm{MNOP}}
\providecommand{\BKM}{\mathrm{BKM}}
\providecommand{\CoHA}{\mathrm{CoHA}}
\providecommand{\red}{\mathrm{red}}
\providecommand{\spin}{\mathrm{spin}}
\providecommand{\un}{\mathrm{un}}
\providecommand{\OP}{\mathrm{OP}}
\providecommand{\reg}{\mathrm{reg}}
\providecommand{\cO}{\mathcal{O}}
\providecommand{\cF}{\mathcal{F}}
\providecommand{\cL}{\mathcal{L}}
\providecommand{\cR}{\mathcal{R}}
\providecommand{\cZ}{\mathcal{Z}}
\providecommand{\fg}{\mathfrak{g}}
\providecommand{\HH}{\mathbb{H}}
\providecommand{\CC}{\mathbb{C}}
\providecommand{\QQ}{\mathbb{Q}}
\providecommand{\ZZ}{\mathbb{Z}}

\section*{Central Regularisation of the Igusa Reciprocal}
\label{sec:central-regularisation-igusa-reciprocal}

The symbol \(L(5/2,\Phi_{10}^{-1})\) is ambiguous until four choices
are fixed: the automorphic form normalisation, the
analytic continuation of the \(L\)-factor, the regularisation of the
reciprocal Siegel form, and the comparison map from the reciprocal
partition function to an automorphic \(L\)-factor.  These choices are
independent.

\subsection*{Igusa and Oberdieck-Pandharipande Normalisations}

Let \(\Delta_5\) denote the theta-normalised Gritsenko-Nikulin form
of weight \(5\) and character \(\nu_{\Delta_5}\), with first Fourier
coefficient
\[
        [q^{1/2}r^{1/2}s^{1/2}]\Delta_5 = 64 .
\]
The monic Borcherds product and the two square conventions are
\[
        D_5=64^{-1}\Delta_5,\qquad
        \Phi_{10}^{\un}=\Delta_5^2,\qquad
        \Phi_{10}^{\OP}=D_5^2=4096^{-1}\Delta_5^2 .
\]
The reduced Oberdieck-Pandharipande scalar for \(K3\times E\) is
\[
        Z_{K3\times E}^{\DT,\red}
        =-\bigl(\Phi_{10}^{\OP}\bigr)^{-1}
        =-D_5^{-2}
        =-4096\,\Delta_5^{-2}.
\]
The sign belongs to the Behrend orientation convention.  The factor
\(4096=64^2\) is a scalar normalisation.  Root multiplicities,
periods, and \(\kappa_\bullet\)-invariants require separate
constructions.

\subsection*{Five Objects}

\[
\begin{array}{lll}
\hline
\text{object} & \text{definition} & \text{analytic behaviour}\\
\hline
\Phi_{10}^{\un} & \Delta_5^2
& \text{holomorphic Siegel cusp form, full level}\\
L_{\spin}(s,\Phi_{10}^{\un}) & \text{Andrianov spinor }L\text{-factor}
& \text{meromorphic continuation by Saito-Kurokawa}\\
(\Phi_{10}^{\OP})^{-1} & D_5^{-2}
& \text{meromorphic Siegel form with divisor poles}\\
Z_{K3\times E}^{\DT,\red} & -(\Phi_{10}^{\OP})^{-1}
& \text{reduced DT partition function}\\
\operatorname{den}(\fg_{\Delta_5}) & D_5(2Z)
& \text{Borcherds-Kac-Moody denominator}\\
\hline
\end{array}
\]

The first line determines a Hecke eigenline.  The third line is a
reciprocal modular form.  The fourth line is an enumerative partition
function.  The fifth line is a denominator identity for
\(\fg_{\Delta_5}\).  A physical amplitude requires, in addition, a
choice of integration domain, measure, infrared subtraction, and
coupling normalisation.

\subsection*{The Spinor \(L\)-Factor}

Put
\[
        \Delta(q)=q\prod_{n\geq1}(1-q^n)^{24},\qquad
        E_6(q)=1-504\sum_{n\geq1}\sigma_5(n)q^n,
\]
\[
        g:=\Delta E_6=\sum_{n\geq1}a_n(g)q^n\in S_{18}(\SL_2(\ZZ)).
\]
The Igusa square \(\Phi_{10}^{\un}\) spans the Saito-Kurokawa
Spezialschar line attached to \(g\).  Andrianov's spinor factor is
\[
        L_{\spin}(s,\Phi_{10}^{\un})
        =\zeta(s-8)\zeta(s-9)L(s,g),
\]
and the completed factor in the local normalisation used here is
\[
        L_{\spin}^{*}(s,\Phi_{10}^{\un})
        =(2\pi)^{-2s}\Gamma(s)\Gamma(s-8)
        L_{\spin}(s,\Phi_{10}^{\un}).
\]
Thus \(L_{\spin}(s,\Phi_{10}^{\un})\) has a simple pole at \(s=10\):
\[
        \operatorname*{Res}_{s=10}
        L_{\spin}(s,\Phi_{10}^{\un})
        =\zeta(2)L(10,g)
        =\frac{\pi^2}{6}L(10,g),
\]
and
\[
        \operatorname*{Res}_{s=10}
        L_{\spin}^{*}(s,\Phi_{10}^{\un})
        =(2\pi)^{-20}\Gamma(10)\Gamma(2)
        \frac{\pi^2}{6}L(10,g).
\]
At \(s=9\), the pole of \(\zeta(s-8)\) is cancelled by \(L(9,g)=0\),
because the level-one weight-\(18\) form \(g\) has root number
\((-1)^9=-1\).

Manin-Deligne periods give
\[
        L(10,g)\in
        \QQ^{\times}(2\pi i)^{10}\Omega^+(g).
\]
This is the strongest period statement available from the standard
Saito-Kurokawa and Eichler-Shimura input alone.  The Fourier coefficient
\[
        a_{10}(g)=541648800
\]
is a coefficient of \(g\).  The value \(L(10,g)\) remains undetermined
until a modular-symbol or period-polynomial computation fixes the
corresponding period ratio.

\subsection*{Programme Coordinate}

If the programme coordinate \(s_{\mathrm{pr}}\) is compared with the
Andrianov coordinate by matching involutions
\[
        s_{\mathrm{pr}}\longmapsto 5-s_{\mathrm{pr}},\qquad
        s\longmapsto 20-s,
\]
the coordinate map is
\[
        \iota(s_{\mathrm{pr}})=4s_{\mathrm{pr}}.
\]
Hence \(s_{\mathrm{pr}}=5/2\) maps to the Andrianov pole \(s=10\).
For the pulled-back spinor factor
\[
        \iota^*L_{\spin}(s_{\mathrm{pr}})
        :=L_{\spin}(4s_{\mathrm{pr}},\Phi_{10}^{\un}),
\]
one has
\[
        \operatorname*{Res}_{s_{\mathrm{pr}}=5/2}
        \iota^*L_{\spin}
        =\frac14\cdot\frac{\pi^2}{6}L(10,g).
\]
For the reciprocal Euler product \(L_{\spin}(4s_{\mathrm{pr}},
\Phi_{10}^{\un})^{-1}\), the same point is a simple zero.  This
reciprocal Euler product is distinct from the Mellin transform of
\((\Phi_{10}^{\OP})^{-1}\).

\subsection*{Regularised Reciprocal of \(\Phi_{10}^{\OP}\)}

Let \(\mathcal F_T\subset \Sp_4(\ZZ)\backslash\HH_2\) be a standard
Siegel fundamental domain truncated at height \(T\).  Let
\(H=H_1\cup H_4\) be the Humbert divisor of \(\Phi_{10}^{\OP}\).  Fix
a kernel \(E(Z,s)\), for instance a degree-two Eisenstein kernel, and
fix tubular neighbourhoods \(N_\epsilon(H)\).  For \(\Re(s)\) large set
\[
        I_{T,\epsilon}(s)
        =
        \int_{\mathcal F_T\setminus N_\epsilon(H)}
        -(\Phi_{10}^{\OP}(Z))^{-1}E(Z,s)\,d\mu(Z).
\]
Let \(P_{T,\epsilon}(s)\) be the explicit polar subtraction determined
by the Fourier expansion at the cusp and by the principal parts along
\(H_1\) and \(H_4\).  The regularised reciprocal value is the finite
part
\[
        I_{\OP}^{\reg}(s;\,E,\mathcal F,P)
        =
        \operatorname*{FP}_{\epsilon=0}
        \operatorname*{FP}_{T=\infty}
        \bigl(I_{T,\epsilon}(s)-P_{T,\epsilon}(s)\bigr).
\]
This is a definition after the data \((E,\mathcal F,P)\) have been
fixed.  A theorem comparing it with the spinor factor must prove an
unfolding formula of the shape
\[
        I_{\OP}^{\reg}(s_{\mathrm{pr}};\,E,\mathcal F,P)
        =
        J_{E,P}(s_{\mathrm{pr}})
        L_{\spin}(4s_{\mathrm{pr}},\Phi_{10}^{\un})
        +R_{E,P}(s_{\mathrm{pr}}),
\]
with \(J_{E,P}\) and \(R_{E,P}\) computed from the same truncation and
principal parts.  Without this unfolding formula, the central
regularised OP/DT value at \(s_{\mathrm{pr}}=5/2\) remains an open
calculation.  The spinor residue above is a necessary input for such a
comparison and still leaves the comparison theorem to prove.

\subsection*{The \(K3\times E\) Scalars}

For \(Y=K3\times E\), the subscripted invariants are
\[
\begin{array}{c|c|l}
\hline
\text{symbol} & \text{value} & \text{construction}\\
\hline
\kcat(Y) & 0 &
\chi(\cO_Y)=\chi(\cO_{K3})\chi(\cO_E)=2\cdot0\\
\kch(Y) & 0 &
\sum_q(-1)^qh^{0,q}(Y)=1-1+1-1\\
\kHeis(Y) & 3 &
\text{Mukai-Heisenberg specialisation}\\
\kBKM(\Delta_5) & 5 &
c_1(0)/2=10/2\\
\kfib(Y) & 24 &
\operatorname{rk}\widetilde H(K3,\ZZ)=1+22+1\\
\hline
\end{array}
\]
The expression \(\kch(Y)+\chi(\cO_{\mathrm{fiber}})\) gives \(0\) for
an elliptic fibre and \(2\) for a K3 fibre.  It therefore cannot
produce \(\kBKM(\Delta_5)=5\).  The value \(5\) is the Borcherds weight
of the denominator form, fixed by \(c_1(0)/2\).

For the primitive CHL denominator ladder
\[
        N\in\{1,2,3,4,6\},
\]
the constants are
\[
\begin{array}{c|ccccc}
N & 1&2&3&4&6\\
\hline
c_N(0) & 10&8&6&4&2\\
\kBKM(\Phi_N)=c_N(0)/2 & 5&4&3&2&1
\end{array}
\]
These are denominator constants.  They are separate from DT/GW
partition functions and from Jatkar-Sen physical square conventions.

\subsection*{Branch Separation}

\begin{enumerate}
\item The \(\BKM\) denominator branch is the Borcherds product
\(\operatorname{den}(\fg_{\Delta_5})=D_5(2Z)\).  A Hall-Drinfeld
double additionally requires a compact critical Hall construction,
negative half, Hopf pairing, coproduct, completion, and associator.
\item The \(\DT/\GW/\MNOP\) branch compares curve-counting partition
functions after the substitution \(-q=e^{iu}\), the degree-zero DT
factor removal, and the relevant wall crossing.  It supplies the
OP/DT scalar above.
\item The class-\(\mathcal S\) and VOA branches use their own
chiral-algebra and conformal-block constructions.  Their central
charges and periods are comparison data only after a named functorial
map has been built.
\item The Yangian/Mukai branch is the K3 Mukai-lattice current branch.
Its Lie-level abelianity and its vertex or double-level interactions
are separate from the \(\BKM\) denominator of \(\fg_{\Delta_5}\).
\item For \(d\geq3\), the native \(\Phi\)-output is \(E_1\)-chiral.
The \(E_2\)-structure lives on a centre, double, or module layer such
as \(\cZ(\operatorname{Rep}^{E_1}(A))\).  In the local
\(\CC^3\) model the \(\CoHA\) gives the positive half \(Y^+\).
The algebra \(\mathcal W_{1+\infty}\) belongs to the double or centre
layer.
\end{enumerate}

\subsection*{Mathematical Conclusion}

The automorphic spinor calculation gives a precise meromorphic fact:
\[
        \operatorname*{Res}_{s=10}
        L_{\spin}(s,\Phi_{10}^{\un})
        =
        \frac{\pi^2}{6}L(10,\Delta E_6),
\]
with
\[
        L(10,\Delta E_6)\in
        \QQ^{\times}(2\pi i)^{10}\Omega^+(\Delta E_6).
\]
Under the coordinate comparison \(s=4s_{\mathrm{pr}}\), the pulled-back
spinor residue at \(s_{\mathrm{pr}}=5/2\) is one quarter of this
residue.

The regularised central value of the reciprocal OP/DT scalar
\[
        -(\Phi_{10}^{\OP})^{-1}
\]
is a different object.  It is determined only after the kernel,
truncation, Humbert-divisor subtraction, and unfolding comparison have
been specified and computed.  The sources cited here leave that
regularised value uncomputed.  The available theorem is
the Saito-Kurokawa spinor factorisation and its residue, together with
the separation of the OP/DT scalar, the BKM denominator, the
Hall-Drinfeld problem, and the physical amplitude.

\subsection*{Primary Anchors}

\begin{itemize}
\item Andrianov 1974: spinor \(L\)-factor of genus-two Siegel cusp
forms and the Saito-Kurokawa factorisation.
\item Maass 1979, Andrianov 1979, Zagier 1981, Evdokimov 1980:
\(\Phi_{10}^{\un}\) as the Saito-Kurokawa lift of
\(g=\Delta E_6\), with
\(L_{\spin}(s)=\zeta(s-8)\zeta(s-9)L(s,g)\).
\item Manin 1972 and Deligne 1979: period normalisation of the
critical values \(L(m,g)\), \(1\leq m\leq17\).
\item Gritsenko-Nikulin 1995 and Borcherds 1998:
\(\Delta_5=\operatorname{Borch}(\phi_{0,1})\),
\(D_5=64^{-1}\Delta_5\), and
\(\kBKM(\Delta_5)=c_1(0)/2=5\).
\item Oberdieck-Pandharipande: the reduced \(K3\times E\) scalar
\(Z_{K3\times E}^{\DT,\red}=-(\Phi_{10}^{\OP})^{-1}\).
\end{itemize}
